% !PRINTED 65
\weberpage{65}
%<<CONT

is obtained, so that here the sign \(\Pi(n)\) has the same meaning as in the first section (\S~7).

Any arrangement of the system \(\mathfrak A\) will be denoted by \(\mathfrak A'\), or more fully, if \(\alpha_1,\alpha_2,\ldots,\alpha_n\) are the digits \(1,2,\ldots,n\) in some order, by
\begin{equation}
 \mathfrak A'=\alpha_1,\alpha_2,\ldots,\alpha_n.
\tag{3}
\end{equation}

One can derive any arrangement from any other in very different ways, that is, one can carry out a permutation in different ways. Among the several possibilities, those carried out by so-called \gesperrt{transpositions}, that is, by successive interchange of only two elements, are of special interest for us. By several transpositions performed successively one can derive every arrangement, for example \(\mathfrak A'\), from \(\mathfrak A\). For this purpose one may proceed, for instance, by first interchanging in \(\mathfrak A\) the element \(1\) with that which stands in the first place in \(\mathfrak A'\), that is, with \(\alpha_1\) (unless \(\alpha_1=1\)); then, if \(\alpha_2\) is not already \(=2\), interchanging \(2\) with \(\alpha_2\), and so on.

Thus, for example, to pass from \((1,2,3,4)\) to \((4,3,2,1)\), one forms the arrangements
\[
 (1,2,3,4),\qquad (4,2,3,1),\qquad (4,3,2,1).
\]
If we briefly denote a transposition by the two interchanged digits, and therefore the interchange of \(1\) with \(2\) by \((1,2)\), then here we have performed successively the transpositions
\[
 (1,4),\quad (2,3).
\]

It is to be observed that the passage from one arrangement to a definite other one can be reached in infinitely many different ways by successive transpositions. Thus the arrangement \((1,2,3,4)\) also passes into \((4,3,2,1)\) by the transpositions
\[
 (1,2),\quad (1,3),\quad (2,4),\quad (1,2).
\]

\begin{center}
{\large \S.\ 19.}

\vspace{0.5em}
{\large\gesperrt{Permutations of the First and Second Kind.}}
\end{center}

The \(\Pi(n)\) arrangements of \(n\) elements can be separated into two kinds according to the following point of view.

From the \(n\) elements of our system one can form
\[
 {n(n-1)\over 2}
\]
and no more pairs. We now assign to the \(n\) elements
%CONT>>
